\section{An adaptive normal-surface fallback in an isolated worker}
\label{sec:normal-surface}

\subsection{A complete geometric decision with a different bottleneck}

The built-in complete backend computes Khovanov homology.  We now add an
optional alternative based on the topology of the knot exterior.  For a
classical knot $K\subset S^3$, its exterior is a solid torus exactly when
$K$ is the unknot.  Thus a complete solid-torus recognition algorithm on
an exact triangulation of the exterior decides both outcomes; a failed
attempt to find a compressing disc is not by itself a negative verdict.

Normal-surface theory provides an algorithmic framework for that decision.
A normal surface in a triangulation with $t$ tetrahedra has coordinates for
four triangle types and three quadrilateral types per tetrahedron.  The
coordinates are nonnegative integers satisfying face-matching equations;
at most one quadrilateral type is permitted in each tetrahedron.  These
constraints encode properly embedded surfaces, including candidates for
compressing discs.  The local quadrilateral alternatives are one source
of exponential search even though the matching equations are linear.
Successful recognition implementations combine such searches with
simplification and crushing operations.

Burton and Ozlen \cite{burton-ozlen} develop a complete branching algorithm
using normal-surface methods and report experimental polynomial behaviour,
while retaining exponential worst-case complexity.  This is motivation for
trying a different complete algorithm when Khovanov work is expensive,
not evidence of a polynomial or subexponential theorem for our inputs.

We use Regina's documented \code{isSolidTorus} decision routine
\cite{regina-solid-torus}.  The wrapper first applies the library's fast
link simplification, constructs an ideal triangulation of the complement,
and calls that routine.  It never substitutes \code{knowsSolidTorus}, whose
Boolean value says whether the property is already known or cheap to
compute, rather than whether the manifold is a solid torus.  The
implementation trusts Regina's conversion, simplification and exact
recognition.  It does not yet export and independently check a sequence of
normal surfaces or crushing operations.

\subsection{Convention checks and the trust boundary}

Our PD row $(a,b,c,d)$ lists ports counterclockwise, with under-strand ports
zero and two.  Regina's PD input starts at the incoming underpass.  The
worker obtains that port $u\in\{0,2\}$ from the validated orientation and
rotates each row by $u$.  Rotation by two preserves the crossing; rotation
by one would interchange over and under and is forbidden.  Normalized arc
labels increase by one for the external API.  The crossing-free input is
explicitly one trivial link component, rather than an empty link.

The worker independently revalidates the supplied PD with
\code{Diagram.from\_pd}.  It checks that the imported link has one
component, is classical and has the same writhe.  After simplification and
complement construction, it checks that the triangulation is valid,
orientable and has exactly one boundary component.  These checks guard
conversion mistakes but are not a replacement for the imported topology.

Every request and result carries a SHA-256 digest of the normalized PD.
The digest binds the response to its input; it is not a mathematical
certificate.  Evidence includes the Regina engine version and installed
distribution version (respectively 7.4 and 7.4.1 in this experiment), initial
and simplified crossing counts, tetrahedron count, elapsed time and an
explicit external-engine trust label.  A result with a wrong digest,
invalid schema, unsupported status, nonzero worker exit or unparsable output
is inconclusive.  Only a completed, well-formed \code{UNKNOT} or
\code{KNOTTED} response can decide the query.  Serialized results involving
such a verdict, including inside connected-sum factors or restarted
queries, explicitly exclude the external wrapper from the built-in
Khovanov running-time statement.

\subsection{Switching after cheap filters, with bounded discarded work}

The optional \code{--regina} flag enables the probe only for an undecided
diagram with at least 32 crossings.  It runs after the existing one-sided
invariant filters and before RIII search and Khovanov construction.  Small
inputs avoid interpreter and native-library startup.  Visible factors are
processed as their own diagrams: an external query never receives the
whole source braid as if it presented a proper factor.

The local allowance is two seconds by default and is configurable with
\code{--regina-seconds}.  Python also permits \code{regina\_seconds=None}
for an uncapped local query.  The standalone \code{regina\_decide} routine
has the same local allowance semantics and accepts diagrams of any size.
Installing \code{fastunknot[normal]} supplies the optional dependency;
the default package remains dependency-free and the default recognizer
does not dispatch to Regina.

The parent launches a fresh worker using the current Python interpreter.
It never imports the native Regina module itself.  Native import, request
conversion, process startup, simplification and recognition all consume
the local query interval.  The parent polls communication in slices of at
most 50 milliseconds, checking the global callback before each slice and
before accepting a result.  Retried communication retains partially sent
input and already received output.  Every exit path kills a live child,
waits for it to finish, and closes all three pipes; this includes local
expiry, global cancellation and exceptions.  Process creation and cleanup
remain operating-system operations, so this is not a hard real-time
latency guarantee.

A missing dependency, local expiry or failed worker records an inconclusive
probe and leaves the original diagram available to the existing RIII and
Khovanov stages.  A global deadline propagates to \code{UNKNOWN} after
cleanup.  Time spent on unsuccessful native work is not reset or refunded.
The \code{max\_objects} parameter still bounds the Khovanov scanner; it is
not a native-memory ceiling.  This worker uses a time allowance rather
than an independent bound on Regina's allocations.

Regina's simplifier makes random choices.  Each worker resets its random
engine to the library default, making these calls reproducible on the same
build and platform.  This does not promise identical simplification paths
across machines or versions, and it does not make correctness probabilistic:
the simplification moves preserve the knot and the terminal recognizer is
an exact decision routine.  Exhaustion remains inconclusive.

This is an adaptive portfolio based on failure of the cheaper stages,
with a local cap before the built-in complete fallback.  It is not
simultaneous racing or a claim of constant-factor optimality relative to
the best possible backend.  A useful next step is to extract independently
checkable geometric witnesses.  No general subexponential or
quasi-polynomial bound follows from the current integration.

\subsection{Measurements with explicit censoring}

The new fixture directory stores PD data generated by Regina 7.4.1 for its
\code{ExampleLink.monster}, \code{gordian} and \code{gst} examples
\cite{regina-examples}.  Their crossing counts are 10, 141 and 48.
The last is the Gompf--Scharlemann--Thompson knot from their slice-ribbon
investigation, not an unknot fixture.  The Conway diagram is also included
in the benchmark from the existing corpus.

\code{normal\_research/benchmark.py} runs five measured shuffled rounds
plus an excluded warm-up for each of four inputs.  Its three arms are the
default pipeline, an identical default control, and the optional portfolio.
Each arm receives a fresh validated PD, a four-second global allowance
and a 50,000-object Khovanov ceiling.  Native attempts always start a new
interpreter; the portfolio's local native allowance is two seconds.

\begin{center}
\begin{tabular}{lrrr}
\toprule
Input & Default median (s) & Portfolio median (s) & Completed default / portfolio\\
\midrule
Monster unknot & 0.000965 & 0.000968 & 5/5\quad/\quad5/5\\
GST knot & 0.001993 & 0.001983 & 5/5\quad/\quad5/5\\
Gordian unknot & $4.043$ (censored) & $1.788$ & 0/5\quad/\quad5/5\\
Conway knot & 0.001041 & 0.001043 & 5/5\quad/\quad5/5\\
\bottomrule
\end{tabular}
\end{center}

All conclusive verdicts agree.  Both default arms exhaust their time
allowance on every measured Gordian query, whereas the portfolio recognizes
it in all five.  The number $4.043$ includes validation and is an observed
resource-limited duration, \emph{not} a completion time or a denominator
for a reported completion-speedup factor.  Native evidence records
simplification from 141 to 111 crossings followed by recognition of a
119-tetrahedron complement.

The dispatch policy matters.  Standalone native queries, with a three-second
allowance, take about 0.278 seconds for Monster and 1.402 seconds for Gordian;
GST exhausts that allowance.  The ordinary Alexander filter rejects GST
in about two milliseconds, so the integrated portfolio avoids the slower
native call.  Monster and Conway remain below the crossing threshold.
These are four concrete examples and one successful hard-unknot improvement,
not a representative average-case study or an asymptotic conclusion.
Raw timing samples and verdict evidence are stored in
\code{results/normal\_surface\_20261008.json}.

Eight new tests check per-crossing conversion, communication across repeated
polls, missing dependencies, malformed responses, input binding, pipeline
dispatch, fallback and actual worker cleanup on local and global cancellation.
With Regina installed they also check six exact native verdicts, global
exhaustion and an actual CLI recognition of the Gordian unknot.  The full
integration log and the dependency-free focused log are retained alongside
the article.  The optional tests are skipped explicitly when Regina is
not installed; the failure-handling tests still run in that configuration.
